% Protocol diagram only: no invented accuracy or learning-curve data.
% Each evaluation uses a separate copy; new learning cannot alter old retests.
\begin{tikzpicture}[
  x=1mm,y=1mm,font=\figurefont\small,
  box/.style={rectangle,draw=BookInk,fill=BookPaper,line width=0.7pt,
    align=center,inner xsep=2mm,inner ysep=2mm,minimum height=12mm},
  flow/.style={draw=BookInk,line width=0.8pt,-{Latex[length=2mm]},
    rounded corners=1mm}
]
  \node[box,text width=60mm,fill=FigureModel!10] (original) at (40,93)
    {原版本\\冻结评测起点，复制两份};
  \node[box,text width=60mm,fill=FigureOutput!10] (updated) at (124,93)
    {更新版本\\冻结评测起点，复制两份};
  \node[box,text width=29mm] (old-a) at (20,62)
    {旧任务复测\\当前有效要求};
  \node[box,text width=29mm] (new-a) at (60,62)
    {同一新任务学习\\等样本、等预算};
  \node[box,text width=29mm] (old-b) at (104,62)
    {旧任务复测\\当前有效要求};
  \node[box,text width=29mm] (new-b) at (144,62)
    {同一新任务学习\\等样本、等预算};
  \draw[flow] (original.south) -- (40,79) -| (old-a.north);
  \draw[flow] (40,79) -| (new-a.north);
  \draw[flow] (updated.south) -- (124,79) -| (old-b.north);
  \draw[flow] (124,79) -| (new-b.north);
  \node[box,text width=29mm,fill=FigureInput!10] (old-ra) at (20,33)
    {原版本的\\旧任务结果};
  \node[box,text width=29mm,fill=FigureLoss!10] (new-ra) at (60,33)
    {原版本的\\后续学习过程};
  \node[box,text width=29mm,fill=FigureInput!10] (old-rb) at (104,33)
    {更新版本的\\旧任务结果};
  \node[box,text width=29mm,fill=FigureLoss!10] (new-rb) at (144,33)
    {更新版本的\\后续学习过程};
  \draw[flow] (old-a) -- (old-ra);
  \draw[flow] (new-a) -- (new-ra);
  \draw[flow] (old-b) -- (old-rb);
  \draw[flow] (new-b) -- (new-rb);
  \node[align=center,text=BookInk,text width=160mm] at (82,6)
    {两侧旧任务结果对照：能力保留；两侧同预算学习对照：可塑性\\
     进一步控制相关历史经验的有无：检验知识积累};
\end{tikzpicture}
